%% kapa: exam=1 round=37 year=2026 period=1 form=A subject=economics subject_ko=경제학원론 seq=4 number=44
%% source: Q-Net 감정평가사 기출문제 (cst00309) artlSeq=5258924 file=raw/1008/5258924/2026년 제37회 감정평가사 1차 1교시.pdf
%% conversion: kordoc 4.24.51 -> tools/gen1.py (auto); review=proofread-sonnet; flags=none
\begin{question}{44}
소비자이론에 관한 설명으로 옳은 것은?
\begin{choices}
  \item 보상수요곡선은 대체효과만을 고려하여 도출한다.
  \item 기펜재(Giffen goods)는 가격이 하락할 때 소비가 증가한다.
  \item 수요의 법칙에 따르면 소득 증가 시 수요는 항상 증가한다.
  \item 두 재화의 한계대체율이 체감하려면 각 재화의 한계효용은 모두 체감해야 한다.
  \item 가격 상승 시 소득효과는 모든 재화의 수요량을 증가시킨다.
\end{choices}
\end{question}
\answer{1}
